\section{Постановка задачи}
\label{sec:Chapter1} \index{Chapter1}

%Необходимо формально изложить суть задачи в данной секции, предоставив такие
%ясные и точные описания, которые позволят в последующем оценить, насколько
%разработанное решение соответствует поставленной задаче. Текст главы должен
%следовать структуре технического задания, включая как описание самой задачи,
%так и набор требований к ее решению.

Целью работы является разработка механизма репликации таблиц трансляций и данных пользовательских приложений, оценка его влияния на производительность в условиях NUMA-нагрузок, а также автоматизация механизма.

Для достижения поставленной цели необходимо решить следующие задачи:
\begin{enumerate}
    \item Спроектировать механизм репликации таблиц трансляций и данных для NUMA-систем, определить его архитектуру и пользовательский интерфейс, а также реализовать механизм в ядре Linux 6.6.

    При проектировании необходимо обеспечить прозрачность механизма для пользовательских приложений и сохранение корректности их выполнения.

    Технические требования к реализации:
    \begin{itemize}
        \item Реализовать в ядре дистрибутива openEuler версии 25.03\cite{openEuler}, основанного на LTS-версии 6.6 ядра Linux и дополненного собственными модификациями и оптимизациями.

        Поддержка данного ядра ведётся в рамках отдельной ветки, синхронизируемой с upstream-проектом Linux. Но отличия между оригинальным Linux 6.6 и openEuler 25.03 в затрагиваемых механизмом подсистемах совсем несущественные, поэтому для большей ясности будет упоминаться именно Linux.
        \item Реализовать под архитектуру arm64.
        \item Поддержать \textit{большие страницы} данных~--- функциональности \textit{Transparent Huge Pages}\cite{THP} и \textit{HugeTLB}\cite{hugetlb}.
        \item Поддержать 4K и 64K размеры страниц.
    \end{itemize}
    \item Выполнить верификацию разработанного механизма и подтвердить корректность его функционирования.
    \item Провести экспериментальную оценку производительности механизма на наборе тестовых приложений.

    Интерес представляют в первую очередь многопоточные тесты, занимающие все NUMA-ноды и взаимодействующие с большим объемом данных. Например, серверы баз данных и фреймворки для обработки больших данных.

    \item Исследовать влияние механизма на производительность приложений, выявить факторы, определяющие эффективность репликации, и на основе полученных результатов построить модель автоматического принятия решений о включении и отключении механизма для конкретных процессов.
\end{enumerate}
